%=============================================================================
%  TECHNIQUES DE HAUTE TENSION — High Voltage Engineering
%  Document pédagogique complet  |  Niveau Licence / Master
%=============================================================================
\documentclass[12pt, a4paper, twoside]{report}

%─── Language & XeLaTeX fonts ────────────────────────────────────────────────
\usepackage[french]{babel}
\usepackage{fontspec}
\defaultfontfeatures{Ligatures=TeX,Scale=MatchLowercase}
\setmainfont{Latin Modern Roman}
\setsansfont{Latin Modern Sans}
\setmonofont{Latin Modern Mono}

%─── Page geometry ────────────────────────────────────────────────────────────
\usepackage[
  a4paper,
  top=2.5cm, bottom=2.8cm,
  left=2.8cm, right=2.2cm,
  headheight=15pt
]{geometry}

%─── Colors ───────────────────────────────────────────────────────────────────
\usepackage[dvipsnames, svgnames, table]{xcolor}
\definecolor{darkblue}{HTML}{0D2B55}
\definecolor{midblue}{HTML}{1A4A8A}
\definecolor{lightblue}{HTML}{3A7CC1}
\definecolor{accent}{HTML}{E8A020}
\definecolor{lightgray}{HTML}{F4F6F9}
\definecolor{bordergray}{HTML}{CBD5E0}
\definecolor{hvred}{HTML}{C0392B}
\definecolor{hvgreen}{HTML}{1A7A4A}
\definecolor{boxbg}{HTML}{EBF4FF}
\definecolor{warnbg}{HTML}{FFF8E7}
\definecolor{warnborder}{HTML}{F6AD55}
\definecolor{successbg}{HTML}{F0FFF4}
\definecolor{successborder}{HTML}{48BB78}
\definecolor{formulabg}{HTML}{EEF2FF}
\definecolor{textdark}{HTML}{1A202C}
\definecolor{textmed}{HTML}{4A5568}

%─── Graphics & Figures ───────────────────────────────────────────────────────
\usepackage{graphicx}
\usepackage{float}
\usepackage{subcaption}
\usepackage{wrapfig}
\graphicspath{{./figures/}}

%─── TikZ & PGFPlots ──────────────────────────────────────────────────────────
\usepackage{tikz}
\usepackage{pgfplots}
\pgfplotsset{compat=1.18}
\usetikzlibrary{
  arrows.meta, shapes.geometric, shapes.misc,
  positioning, calc, decorations.pathmorphing,
  decorations.markings, patterns, shadows,
  backgrounds, fit, matrix, circuits.ee.IEC
}

%─── Math ─────────────────────────────────────────────────────────────────────
\usepackage{amsmath, amssymb, amsthm, mathtools}
\usepackage{siunitx}
\sisetup{locale=FR, per-mode=fraction}

%─── Tables ───────────────────────────────────────────────────────────────────
\usepackage{booktabs}
\usepackage{tabularx}
\usepackage{multirow}
\usepackage{array}
\usepackage{longtable}
\usepackage{colortbl}

%─── Typography & Layout ──────────────────────────────────────────────────────
\usepackage{microtype}
\usepackage{setspace}
\usepackage{parskip}
\setlength{\parskip}{6pt}
\setlength{\parindent}{0pt}

%─── Custom boxes ─────────────────────────────────────────────────────────────
\usepackage[most]{tcolorbox}
\tcbuselibrary{skins, breakable, listingsutf8}

%─── Headers & Footers ────────────────────────────────────────────────────────
\usepackage{fancyhdr}
\pagestyle{fancy}
\fancyhf{}
\fancyhead[LE,RO]{\small\thepage}
\fancyhead[LO]{\small\textcolor{midblue}{\leftmark}}
\fancyhead[RE]{\small\textcolor{midblue}{Techniques de Haute Tension}}
\fancyfoot[C]{\small\textcolor{textmed}{Support de cours — Ingénierie Électrique}}
\renewcommand{\headrulewidth}{0.4pt}
\renewcommand{\footrulewidth}{0.3pt}

%─── Chapter/Section styling ─────────────────────────────────────────────────
\usepackage{titlesec}
\titleformat{\chapter}[block]
{\normalfont\huge\bfseries\color{white}}
{}
{0pt}
{
	\colorbox{darkblue}{
		\parbox{\dimexpr\textwidth-2\fboxsep\relax}{
			\hspace{4pt}\textcolor{accent}{\thechapter.}\hspace{8pt}\strut ##1\hspace{4pt}
		}
	}
}
\titlespacing*{\chapter}{0pt}{0pt}{20pt}

\titleformat{\section}
  {\normalfont\large\bfseries\color{darkblue}}
  {\textcolor{accent}{\thesection}\hspace{8pt}}
  {0pt}{}
  [\vspace{-4pt}\textcolor{lightblue}{\rule{\textwidth}{1pt}}]
\titlespacing*{\section}{0pt}{16pt}{8pt}

\titleformat{\subsection}
  {\normalfont\normalsize\bfseries\color{midblue}}
  {\textcolor{midblue}{\thesubsection}\hspace{6pt}}
  {0pt}{}
\titlespacing*{\subsection}{0pt}{12pt}{4pt}

\titleformat{\subsubsection}
  {\normalfont\normalsize\itshape\color{lightblue}}
  {\thesubsubsection\hspace{6pt}}
  {0pt}{}
\titlespacing*{\subsubsection}{0pt}{8pt}{2pt}

%─── TOC styling ──────────────────────────────────────────────────────────────
% tocloft disabled to avoid \@starttoc redefinition conflict in this setup

%─── Hyperref ─────────────────────────────────────────────────────────────────
\usepackage[
  colorlinks=true,
  linkcolor=midblue,
  citecolor=hvgreen,
  urlcolor=lightblue,
  pdfauthor={HV Engineering Course},
  pdftitle={Techniques de Haute Tension},
  pdfsubject={High Voltage Engineering}
]{hyperref}
\usepackage{bookmark}

%─── Custom environments ─────────────────────────────────────────────────────
% Info box
\newtcolorbox{infobox}[1]{
  enhanced, breakable,
  colback=boxbg, colframe=lightblue, colbacktitle=midblue,
  title={\faIcon{info-circle}\hspace{4pt}#1},
  fonttitle=\bfseries\color{white},
  boxrule=1pt, arc=4pt,
  leftrule=4pt,
  before skip=8pt, after skip=8pt
}
% Warning box
\newtcolorbox{warnbox}[1]{
  enhanced, breakable,
  colback=warnbg, colframe=warnborder, colbacktitle=accent,
  title={\textbf{⚠\space#1}},
  fonttitle=\bfseries\color{white},
  boxrule=1pt, arc=4pt, leftrule=4pt,
  before skip=8pt, after skip=8pt
}
% Success / À retenir box
\newtcolorbox{retenir}[1]{
  enhanced, breakable,
  colback=successbg, colframe=successborder, colbacktitle=hvgreen,
  title={\textbf{✓\space#1}},
  fonttitle=\bfseries\color{white},
  boxrule=1pt, arc=4pt, leftrule=4pt,
  before skip=8pt, after skip=8pt
}
% Formula box
\newtcolorbox{formulebox}[1]{
  enhanced,
  colback=formulabg, colframe=midblue, colbacktitle=darkblue,
  title={\textbf{#1}},
  fonttitle=\bfseries\color{white},
  boxrule=1.2pt, arc=5pt,
  before skip=8pt, after skip=8pt
}
% Example box
\newtcolorbox{exemple}[1]{
  enhanced, breakable,
  colback=white, colframe=accent, colbacktitle=accent,
  title={\textbf{Exemple : #1}},
  fonttitle=\bfseries\color{white},
  boxrule=1pt, arc=4pt, leftrule=4pt,
  before skip=8pt, after skip=8pt
}

%─── Misc ─────────────────────────────────────────────────────────────────────
\usepackage{enumitem}
\setlist[itemize]{leftmargin=1.5em, itemsep=3pt, topsep=4pt}
\setlist[enumerate]{leftmargin=1.5em, itemsep=3pt, topsep=4pt}
\usepackage{caption}
\captionsetup{font=small, labelfont={bf,color=midblue}, labelsep=period}
\usepackage{epigraph}
\usepackage{afterpage}
\usepackage{pdflscape}

%─── Font icons (optional, comment out if not installed) ──────────────────────
% \usepackage{fontawesome5}
\providecommand{\faIcon}[1]{}

%=============================================================================
%  DOCUMENT
%=============================================================================
\begin{document}

%─── COVER PAGE ───────────────────────────────────────────────────────────────
\begin{titlepage}
\pagecolor{darkblue}
\afterpage{\nopagecolor}
\centering
\vspace*{1.5cm}

{\color{accent}\rule{\textwidth}{3pt}}
\vspace{0.4cm}

{\color{white}\fontsize{36}{44}\selectfont\bfseries
  TECHNIQUES DE HAUTE TENSION}

\vspace{0.3cm}
{\color{lightblue}\fontsize{20}{26}\selectfont
  High Voltage Engineering}

\vspace{0.4cm}
{\color{accent}\rule{\textwidth}{1.5pt}}
\vspace{0.6cm}

% Cover figure placeholder
\begin{tikzpicture}
  \fill[darkblue!60] (0,0) rectangle (14,6);
  % Transmission tower sketch
  \draw[accent!80, line width=1.5pt] (2,0) -- (2,5.5);
  \draw[accent!80, line width=1.5pt] (1,0) -- (2,2) -- (3,0);
  \draw[accent!80, line width=1.5pt] (0.5,1.5) -- (3.5,1.5);
  \draw[accent!80, line width=1.5pt] (0.8,3.0) -- (3.2,3.0);
  \draw[accent!80, line width=1.5pt] (1.2,4.5) -- (2.8,4.5);
  % Corona glow
  \foreach \r in {0.3,0.6,0.9,1.2}{
    \draw[lightblue!\the\numexpr120-\r*40\relax, opacity=0.4, line width=\r pt]
      (7,3) circle (\r);
  }
  \fill[white] (7,3) circle (0.18);
  \node[color=white, font=\tiny] at (7,3) {HT};
  % E field arrows
  \foreach \a in {0,45,...,315}{
    \draw[-Stealth, lightblue, opacity=0.6, line width=0.8pt]
      ({7+0.25*cos(\a)},{3+0.25*sin(\a)}) --
      ({7+1.4*cos(\a)},{3+1.4*sin(\a)});
  }
  % Labels
  \node[color=white, font=\large\bfseries] at (7,5.3) {Champ électrique \& Corona};
  \node[color=accent, font=\normalsize] at (2,5.8) {Pylône THT};
  % Wave
  \draw[hvred, line width=1.5pt, domain=9:13.5, samples=100]
    plot(\x, {3 + 2*exp(-(\x-9)/1.2)*sin(720*(\x-9))});
  \node[color=white, font=\normalsize] at (11.2,5.3) {Onde d'impulsion};
  \node[color=hvred, font=\small] at (11.2,0.4) {1.2/50 \(\mu\)s};
\end{tikzpicture}

\vspace{0.6cm}
\begin{tcolorbox}[
  colback=midblue!80, colframe=lightblue, arc=6pt,
  width=0.85\textwidth, boxrule=1pt
]
\centering
\color{white}
\textbf{\large Support de cours — Niveau Licence / Master}\\[4pt]
Ingénierie Électrique \textbullet{} Génie Électrique\\[8pt]
\textcolor{accent}{\rule{0.5\textwidth}{0.5pt}}\\[8pt]
{\small
\begin{tabular}{ll}
  \textbf{Thèmes couverts :} & Champs HT, Claquage, Corona, Paschen\\
  & Surtensions, Génération, Mesure\\
  & Coordination d'isolation, GIS\\
  & Applications industrielles, Normes IEC/IEEE
\end{tabular}
}
\end{tcolorbox}

\vfill
{\color{lightblue!70}\small\textit{Document pédagogique complet avec exemples numériques, graphiques et schémas}}\\[4pt]
{\color{accent}\rule{\textwidth}{1pt}}
{\color{textmed!50}\small\copyright~2024 — HV Engineering Course}
\end{titlepage}

%─── TABLE OF CONTENTS ────────────────────────────────────────────────────────
\tableofcontents
\clearpage

%=============================================================================
\chapter{Introduction}
%=============================================================================

\epigraph{\textit{``La haute tension est à la fois la clé du transport d'énergie moderne et
le plus grand défi de l'ingénieur électricien.''}}{}

\section{Définition des Hautes Tensions}

La \textbf{haute tension} (HT) est définie par la norme \textbf{IEC 60038} comme toute
tension dépassant certains seuils selon la nature du courant :

\begin{center}
\begin{tabular}{lccc}
\toprule
\rowcolor{darkblue}
\textcolor{white}{\textbf{Catégorie}} &
\textcolor{white}{\textbf{Tension CA (rms)}} &
\textcolor{white}{\textbf{Tension CC}} &
\textcolor{white}{\textbf{Symbole}} \\
\midrule
Basse tension (BT)       & $< 1\,\si{kV}$        & $< 1.5\,\si{kV}$  & LV \\
\rowcolor{lightgray}
Haute tension (HT)       & $1 - 52\,\si{kV}$     & $1.5 - 75\,\si{kV}$ & HV \\
Très haute tension (THT) & $52 - 800\,\si{kV}$   & $75 - 750\,\si{kV}$  & EHV \\
\rowcolor{lightgray}
Ultra haute tension (UHT)& $> 800\,\si{kV}$      & $> 750\,\si{kV}$     & UHV \\
\bottomrule
\end{tabular}
\end{center}

\section{Domaines d'Application}

Les hautes tensions sont présentes dans pratiquement tous les secteurs de l'énergie :

\begin{itemize}
  \item \textbf{Transport d'énergie électrique :} Les lignes THT (225\,kV, 400\,kV, 765\,kV)
    minimisent les pertes Joule $P = RI^2$ sur de longues distances en réduisant le courant
    $I = P/({\sqrt{3}\,V})$.

  \item \textbf{Réseaux HVDC :} Liaisons en courant continu haute tension ($\pm$500\,kV à
    $\pm$800\,kV) pour interconnexions continentales et raccordement d'éoliennes offshore.

  \item \textbf{Industrie lourde :} Fours à arc (jusqu'à 100\,kV), électrolyse de l'aluminium
    (700--1000\,V, 300\,kA), précipitateurs électrostatiques (50--100\,kV).

  \item \textbf{Tests et qualification :} Essais de tenue diélectrique des équipements
    (câbles, transformateurs, disjoncteurs, traversées).

  \item \textbf{Applications scientifiques et médicales :} Accélérateurs de particules,
    appareils de radiologie (40--150\,kV), lasers CO$_2$ industriels, spectromètres de masse.

  \item \textbf{Télécommunications :} Alimentation des amplificateurs à tubes, radars,
    émetteurs AM (50--400\,kW, alimentation HT 10--30\,kV).
\end{itemize}

\begin{retenir}{Pourquoi transporter en haute tension ?}
La puissance transportée est $P = \sqrt{3}\,V\,I\,\cos\varphi$.
Pour une puissance donnée, augmenter $V$ réduit $I$ proportionnellement.
Les pertes en ligne étant $P_{\text{pertes}} = 3\,R\,I^2$, une ligne à 400\,kV
transporte la même puissance qu'une ligne 20\,kV avec \textbf{400 fois moins de pertes} !
\end{retenir}

\section{Importance dans les Systèmes Modernes}

La transition énergétique mondiale rend l'ingénierie HT encore plus stratégique :
les énergies renouvelables (éolien offshore, solaire CSP en désert) sont souvent situées
loin des centres de consommation. Le transport HT/HVDC est la seule solution technologique
pour véhiculer ces gigawatts sur des milliers de kilomètres avec des pertes acceptables.

\begin{figure}[H]
\centering
\begin{tikzpicture}[scale=0.85]
  % Simple overview of power system
  \node[draw=hvgreen, fill=hvgreen!15, rounded corners, minimum width=2.2cm, minimum height=1cm, align=center]
    (gen) at (0,0) {\shortstack{\small\textbf{Génération}\\[-2pt]\tiny 10--20\,kV}};
  \node[draw=midblue, fill=midblue!15, rounded corners, minimum width=2.2cm, minimum height=1cm, align=center]
    (step1) at (3.5,0) {\shortstack{\small\textbf{Élévation}\\[-2pt]\tiny 400\,kV}};
  \node[draw=darkblue, fill=darkblue!15, rounded corners, minimum width=2.6cm, minimum height=1cm, text=white, align=center]
    (trans) at (7.5,0) {\shortstack{\small\textbf{Transport THT}\\[-2pt]\tiny 225--765\,kV}};
  \node[draw=accent, fill=accent!15, rounded corners, minimum width=2.2cm, minimum height=1cm, align=center]
    (step2) at (11.5,0) {\shortstack{\small\textbf{Abaissement}\\[-2pt]\tiny 63\,kV}};
  \node[draw=hvred, fill=hvred!10, rounded corners, minimum width=2.2cm, minimum height=1cm, align=center]
    (dist) at (15,0) {\shortstack{\small\textbf{Distribution}\\[-2pt]\tiny 20\,kV/400\,V}};

  \draw[-Stealth, thick, midblue] (gen) -- (step1) node[midway, above, font=\small] {};
  \draw[-Stealth, thick, darkblue] (step1) -- (trans);
  \draw[-Stealth, thick, darkblue] (trans) -- (step2);
  \draw[-Stealth, thick, hvred] (step2) -- (dist);

  % Labels
  \node[below=0.3cm of gen, textmed, font=\footnotesize] {Centrale};
  \node[below=0.3cm of step1, textmed, font=\footnotesize] {Transfo HT};
  \node[below=0.3cm of trans, textmed, font=\footnotesize] {Lignes HT/THT};
  \node[below=0.3cm of step2, textmed, font=\footnotesize] {Transfo BT};
  \node[below=0.3cm of dist, textmed, font=\footnotesize] {Consommateurs};
\end{tikzpicture}
\caption{Schéma simplifié du réseau électrique — du producteur au consommateur}
\end{figure}

%=============================================================================
\chapter{Rappels Fondamentaux}
%=============================================================================

\section{Champ Électrique et Potentiel}

\subsection{Définitions}

Le \textbf{champ électrique} $\vec{E}$ en un point est la force par unité de charge
positive en ce point :
\begin{equation}
  \vec{E} = -\nabla V = -\text{grad}\,V \quad [\si{V/m}]
\end{equation}

Pour une charge ponctuelle $Q$ dans le vide :
\begin{equation}
  E(r) = \frac{Q}{4\pi\varepsilon_0\,r^2} \quad
  \text{avec } \varepsilon_0 = 8.854 \times 10^{-12}\,\si{F/m}
\end{equation}

Le \textbf{potentiel électrique} $V$ en un point $P$ par rapport à l'infini est :
\begin{equation}
  V(P) = -\int_{\infty}^{P} \vec{E}\cdot d\vec{l} = \frac{Q}{4\pi\varepsilon_0\,r}
\end{equation}

\subsection{Champ Uniforme — Plaques Parallèles}

Entre deux plaques conductrices infinies séparées d'une distance $d$ soumises à
une différence de potentiel $U$, le champ est \textbf{uniforme} :
\begin{equation}
  \boxed{E = \frac{U}{d}} \quad [\si{V/m}]
\end{equation}

\begin{exemple}{Calcul du champ entre plaques parallèles}
\textbf{Données :} $U = 100\,\si{kV}$, $d = 10\,\si{cm} = 0.10\,\si{m}$\\[4pt]
\textbf{Calcul :}
\[
  E = \frac{U}{d} = \frac{100\,000}{0.10} = 1\,000\,000\,\si{V/m} = 1\,\si{MV/m}
\]
\textbf{Comparaison :} Le claquage de l'air se produit à $E_{\text{claquage}} \approx 3\,\si{MV/m}$\\
$\Rightarrow$ Marge de sécurité = $3\,\si{MV/m} / 1\,\si{MV/m} = 3$ (facteur 3).
\end{exemple}

\begin{figure}[H]
\centering
\begin{subfigure}{0.45\textwidth}
\centering
\begin{tikzpicture}[scale=0.9]
  % Plates
  \fill[darkblue] (-0.3,-2) rectangle (0,2);
  \fill[hvred] (4,-2) rectangle (4.3,2);
  % Field lines with arrows
  \foreach \y in {-1.5,-0.9,-0.3,0.3,0.9,1.5}{
    \draw[-Stealth, lightblue, line width=1.2pt] (0.1,\y) -- (3.9,\y);
  }
  % Labels
  \node[left, darkblue, font=\large\bfseries] at (-0.3,0) {$-$};
  \node[right, hvred, font=\large\bfseries] at (4.3,0) {$+$};
  \node[below, midblue, font=\small\itshape] at (2,-2.3) {$E = U/d$ (uniforme)};
  \draw[<->, thick, accent] (0.1,-2.5) -- (3.9,-2.5) node[midway,below]{\small $d$};
\end{tikzpicture}
\caption{Champ uniforme — plaques parallèles}
\end{subfigure}
\hfill
\begin{subfigure}{0.45\textwidth}
\centering
\begin{tikzpicture}[scale=0.9]
  % Field lines radiating
  \foreach \angle in {0,30,...,330}{
    \draw[-Stealth, hvred!70, line width=1pt]
      ({0.3*cos(\angle)},{0.3*sin(\angle)}) --
      ({2.2*cos(\angle)},{2.2*sin(\angle)});
  }
  % Equipotential circles
  \foreach \r in {0.8, 1.3, 1.8}{
    \draw[lightblue, dashed, opacity=0.6] (0,0) circle (\r);
  }
  \fill[hvred] (0,0) circle (0.25);
  \node[white, font=\tiny\bfseries] at (0,0) {$+Q$};
  \node[below, midblue, font=\small\itshape] at (0,-2.5)
    {$E = Q/(4\pi\varepsilon_0 r^2)$};
  \node[lightblue, font=\tiny] at (1.3, 0.9) {équipotentielles};
\end{tikzpicture}
\caption{Champ non uniforme — charge ponctuelle}
\end{subfigure}
\caption{Champ électrique uniforme et non uniforme. Les lignes bleues représentent
les lignes de champ, les cercles tiretés les équipotentielles.}
\end{figure}

\section{Loi de Gauss}

La \textbf{loi de Gauss} relie le flux du champ électrique à travers une surface
fermée $\mathcal{S}$ à la charge totale $Q_{\text{enc}}$ enfermée :
\begin{equation}
  \oint_{\mathcal{S}} \vec{E}\cdot d\vec{A} = \frac{Q_{\text{enc}}}{\varepsilon_0}
\end{equation}

\subsection{Application — Câble Coaxial (Câble HT)}

\begin{wrapfigure}{r}{0.38\textwidth}
\vspace{-10pt}
\begin{tikzpicture}[scale=1.0]
  % Outer conductor
  \fill[bordergray!60] (0,0) circle (2.2);
  % Insulation
  \fill[accent!20] (0,0) circle (2.0);
  % Inner conductor
  \fill[midblue!60] (0,0) circle (0.6);
  % Labels
  \draw[<->, thick] (0,0) -- (0.6,0) node[midway, above, font=\small]{$a$};
  \draw[<->, thick] (0,0) -- (0,-2.0) node[midway, right, font=\small]{$R$};
  \node[white, font=\small\bfseries] at (0,0) {HT};
  \node[font=\small, textmed] at (0, -2.5) {Isolant ($\varepsilon_r$)};
  % E field arrow
  \draw[-Stealth, hvred, thick] (0.65,0) -- (1.8,0)
    node[midway, above, font=\small]{$E(r)$};
\end{tikzpicture}
\vspace{-6pt}
\caption{Câble coaxial}
\end{wrapfigure}

Pour un câble coaxial de rayon interne $a$ et externe $R$, la loi de Gauss
en coordonnées cylindriques donne :
\begin{equation}
  \boxed{E(r) = \frac{U}{r\,\ln\!\left(\dfrac{R}{a}\right)}}
  \quad (a \leq r \leq R)
\end{equation}

Le champ est \textbf{maximal} en $r = a$ (surface interne) :
\begin{equation}
  E_{\max} = \frac{U}{a\,\ln\!\left(\dfrac{R}{a}\right)}
\end{equation}

Le rapport optimal $R/a$ minimisant $E_{\max}$ pour un rayon extérieur $R$ fixé
est obtenu en résolvant $\partial E_{\max}/\partial a = 0$, ce qui donne $R/a = e \approx 2.718$.

\begin{infobox}{Implication pour la conception des câbles HT}
Le concepteur cherche à maximiser $a$ (rayon du conducteur) pour réduire $E_{\max}$,
tout en limitant le diamètre extérieur $R$ pour des raisons de coût. Le compromis
optimal donne $R/a \approx e$ pour un câble monophasé.
\end{infobox}

\section{Distribution du Champ — Géométries Courantes}

\begin{center}
\begin{tabularx}{\textwidth}{lXX}
\toprule
\rowcolor{darkblue}
\textcolor{white}{\textbf{Géométrie}} &
\textcolor{white}{\textbf{Expression de $E(r)$ ou $E$}} &
\textcolor{white}{\textbf{Caractéristique}} \\
\midrule
Plaques parallèles & $E = U/d$ & Uniforme \\
\rowcolor{lightgray}
Sphères concentriques & $E(r) = \dfrac{U\,ab}{b-a}\cdot\dfrac{1}{r^2}$ & $\propto 1/r^2$ \\
Cylindres coaxiaux & $E(r) = \dfrac{U}{r\ln(R/a)}$ & $\propto 1/r$ \\
\rowcolor{lightgray}
Pointe--plan & Numérique / FEM & Très non uniforme \\
Sphère--plan & Numérique / FEM & Non uniforme \\
\bottomrule
\end{tabularx}
\end{center}

\section{Propriétés des Matériaux Diélectriques}

\subsection{Permittivité et Polarisation}

Dans un diélectrique, le déplacement électrique est :
\begin{equation}
  \vec{D} = \varepsilon_0\,\varepsilon_r\,\vec{E} = \varepsilon_0\,\vec{E} + \vec{P}
\end{equation}
où $\vec{P} = \varepsilon_0(\varepsilon_r - 1)\vec{E}$ est la \textbf{polarisation diélectrique}.
La polarisation peut être dipolaire (orientationnelle), ionique ou électronique.

\subsection{Tableau des Matériaux Isolants HT}

\begin{center}
\begin{tabularx}{\textwidth}{lccXl}
\toprule
\rowcolor{darkblue}
\textcolor{white}{\textbf{Matériau}} &
\textcolor{white}{\textbf{$\varepsilon_r$}} &
\textcolor{white}{\textbf{$E_{\rm rupt}$ (kV/mm)}} &
\textcolor{white}{\textbf{Avantages}} &
\textcolor{white}{\textbf{Application}} \\
\midrule
Air (1 atm, 20°C) & 1.0006 & $\approx$3 & Gratuit, auto-cicatrisant & Lignes aériennes \\
\rowcolor{lightgray}
SF$_6$ (3 bar, 20°C) & 1.002 & 8--9 & Rigidité 3× l'air & Disjoncteurs GIS \\
Huile minérale & 2.2--2.5 & 10--15 & Bonne tenue thermique & Transformateurs \\
\rowcolor{lightgray}
Papier imprégné huile & 3.5--4.0 & 20--30 & Combiné thermique & Câbles OFHV \\
XLPE (polyéthylène réticulé) & 2.3 & 20--40 & Léger, pas de liquide & Câbles THT modernes \\
\rowcolor{lightgray}
Porcelaine & 6--7 & 10--12 & Résistance mécanique & Isolateurs aériens \\
Verre trempé & 7--8 & 10--14 & Défaut visible & Chaînes d'isolateurs \\
\rowcolor{lightgray}
Silicone RTV & 3.5--4.5 & 15--20 & Hydrophobicité & Isolateurs pollution \\
Époxy (résine) & 3.5--5.0 & 15--25 & Bonne moulabilité & Traversées HT \\
\bottomrule
\end{tabularx}
\end{center}

\begin{warnbox}{Dépendance en fréquence et en température}
La permittivité $\varepsilon_r$ et la rigidité diélectrique varient avec la fréquence, la
température et l'humidité. Une augmentation de température de 20°C peut réduire la rigidité
de l'huile de 20--30\%. Les essais doivent donc être réalisés dans des conditions normalisées
(23°C ± 2°C, 50\% HR selon IEC 60243).
\end{warnbox}

%=============================================================================
\chapter{Phénomènes en Haute Tension}
%=============================================================================

\section{Mécanismes de Claquage dans les Gaz}

\subsection{Théorie de Townsend}

Le \textbf{claquage gazeux} résulte d'une avalanche électronique auto-entretenue.
La théorie de Townsend (1910) repose sur deux coefficients :

\begin{formulebox}{Coefficients de Townsend}
\begin{align}
  &\alpha \quad \text{(1\textsuperscript{er} coefficient)} = \text{ionisations par électron par unité de longueur} \\
  &\gamma \quad \text{(2\textsuperscript{e} coefficient)} = \text{électrons secondaires émis par ion frappant la cathode}
\end{align}

Après $d$ de parcours, un électron génère $e^{\alpha d}$ électrons.
\textbf{Condition de claquage de Townsend :}
\[
  \gamma\left(e^{\alpha d} - 1\right) = 1
  \quad \Longrightarrow \quad
  e^{\alpha d} = 1 + \frac{1}{\gamma}
\]
\end{formulebox}

\begin{exemple}{Calcul du coefficient $\alpha$ dans l'air}
Dans l'air à pression atmosphérique, $\alpha$ est approximé par la formule de Townsend :
\[
  \alpha = A \cdot p \cdot \exp\!\left(-\frac{B\cdot p}{E}\right)
\]
avec pour l'air : $A = 15\,\si{cm^{-1}bar^{-1}}$, $B = 365\,\si{V/(cm\cdot bar)}$\\[6pt]
Pour $E = 30\,\si{kV/cm}$ et $p = 1\,\si{bar}$ :
\[
  \alpha = 15 \times 1 \times \exp\!\left(-\frac{365}{30\,000 \times 10^{-3}}\right)
  \approx 15\,\si{cm^{-1}}
\]
\end{exemple}

\subsection{Mécanisme de Streamer}

Pour les grands espaces et les hautes pressions, le mécanisme de \textbf{streamer} prédomine :
\begin{itemize}
  \item Les charges d'espace créées par l'avalanche déforment le champ local.
  \item Des filaments ionisés se propagent à grande vitesse ($10^5$--$10^6\,\si{m/s}$).
  \item Le claquage se produit en une seule traversée (sans nécessiter $\gamma$).
  \item Condition de streamer (Meek) : $\alpha x_c \approx 18$--20 pour l'air.
\end{itemize}

\section{Claquage dans les Liquides}

Dans les isolants liquides (huile de transformateur), plusieurs mécanismes coexistent :

\begin{itemize}
  \item \textbf{Formation de bulles gazeuses :} cavitation par décharges locales, vaporisation
    par échauffement $\rightarrow$ claquage dans la bulle gazeuse (moins rigide).
  \item \textbf{Migration de particules :} les particules conductrices ou semi-conductrices
    s'alignent sous l'effet du champ et forment un pont conducteur (``bridging'').
  \item \textbf{Ionisation directe :} pour des champs très élevés ($>$100\,MV/m), injection
    directe d'électrons depuis les électrodes.
  \item \textbf{Claquage de zone :} zones de faible rigidité autour des impuretés.
\end{itemize}

\begin{infobox}{Effet de la propreté sur la rigidité de l'huile}
Une huile de transformateur propre (filtrée, déshydratée, dégazée) peut atteindre
$E_{\text{claquage}} \geq 80\,\si{kV/cm}$ (essai normalisé IEC 60156).
La présence de $0.01\%$ d'eau peut la réduire à moins de $20\,\si{kV/cm}$ !
\end{infobox}

\section{Claquage dans les Solides}

Le claquage des solides est généralement \textbf{irréversible} et catastrophique :

\begin{enumerate}
  \item \textbf{Claquage électronique pur} (intrinsèque) : très rapide ($< 10^{-8}\,\si{s}$),
    dû à une avalanche dans la bande de conduction. $E_{\text{claquage}}$ peut atteindre
    $10^3\,\si{MV/m}$ pour des films minces parfaits.
  \item \textbf{Claquage thermique :} la puissance dissipée $P = \sigma E^2$ augmente
    la conductivité $\sigma$ $\rightarrow$ emballement thermique. Prédomine à haute
    température ou fréquence.
  \item \textbf{Claquage électromécanique :} les forces de Maxwell $P_{\text{méca}} = \varepsilon_0\varepsilon_r E^2/2$
    compriment le diélectrique jusqu'à son effondrement mécanique.
  \item \textbf{Claquage par ``treeing'' :} formation d'arborescences de décharges
    partielles progressant sur des heures à des années avant rupture finale.
\end{enumerate}

\section{Loi de Paschen}

\subsection{Énoncé et Formulation}

La \textbf{loi de Paschen} (1889) stipule que la tension de claquage $V_b$ dans un gaz
homogène est uniquement fonction du produit pression $\times$ distance ($p \cdot d$) :
\begin{equation}
  \boxed{V_b = f(p \cdot d)}
\end{equation}

La formulation analytique est :
\begin{equation}
  V_b = \frac{B\,(p\cdot d)}{\ln(A\,p\cdot d) - \ln\!\left[\ln\!\left(1 + \dfrac{1}{\gamma}\right)\right]}
  \label{eq:paschen}
\end{equation}

Pour l'air : $A = 11.5\,\si{Pa^{-1}m^{-1}}$, $B = 364\,\si{V\cdot Pa^{-1}m^{-1}}$,
$\gamma \approx 0.01$.

\subsection{Courbe de Paschen et Minimum}

\begin{figure}[H]
\centering
\begin{tikzpicture}
\begin{semilogxaxis}[
  width=0.85\textwidth, height=7cm,
  xlabel={Produit $p \cdot d$ (Pa$\cdot$m)},
  ylabel={Tension de claquage $V_b$ (V)},
  xmin=0.5, xmax=1000,
  ymin=0, ymax=3000,
  grid=both, grid style={dotted, gray!40},
  legend pos=north east,
  title style={color=darkblue},
  title={Courbe de Paschen — Air ($\gamma = 0.01$)},
  axis background/.style={fill=lightgray!30},
  tick label style={font=\small},
  label style={font=\small, color=textdark},
]
% Paschen curve for air (approximated numerically)
\addplot[
  darkblue, thick, smooth, line width=2pt,
  domain=1:1000, samples=200,
] {364*x/(ln(11.5*x) - ln(ln(1+100)))};

% Minimum marker
\addplot[
  only marks, mark=*, mark size=4pt,
  accent, line width=2pt
] coordinates {(7.5, 327)};

% Annotation lines
\addplot[dashed, accent, opacity=0.7] coordinates {(7.5, 0) (7.5, 327)};
\addplot[dashed, accent, opacity=0.7] coordinates {(0.5, 327) (7.5, 327)};

\node[anchor=west, font=\small, color=hvred] at (axis cs:9, 370)
  {Minimum : $V_{\min}\approx327\,\text{V}$};
\node[anchor=west, font=\small, color=hvred] at (axis cs:9, 220)
  {$p\cdot d \approx 0.76\,\text{Pa\!\cdot\!m}$};

\addlegendentry{$V_b = f(p\cdot d)$ — Air}
\addlegendentry{Minimum de Paschen}

\end{semilogxaxis}
\end{tikzpicture}
\caption{Courbe de Paschen pour l'air. La tension de claquage présente un minimum
(``minimum de Paschen'') autour de $p\cdot d \approx 0.76\,\si{Pa\cdot m}$.}
\end{figure}

\begin{retenir}{Interprétation physique de la courbe de Paschen}
\begin{itemize}
  \item \textbf{À gauche du minimum (faible $p\cdot d$)} : les électrons parcourent de
    longues distances sans collision $\rightarrow$ peu d'ionisations $\rightarrow$ besoin de
    champ très élevé.
  \item \textbf{À droite du minimum (fort $p\cdot d$)} : les électrons perdent leur énergie
    rapidement par collisions $\rightarrow$ besoin de champ plus élevé pour maintenir
    l'ionisation.
  \item \textbf{Au minimum} : conditions optimales d'ionisation. C'est la zone la plus
    ``dangereuse'' car le claquage se produit à la tension minimale.
\end{itemize}
Applications : vide (gauche), pression atmosphérique (droite), SF$_6$ sous pression (très à droite).
\end{retenir}

\section{Effet Corona}

\subsection{Définition et Mécanisme}

L'\textbf{effet corona} (ou décharge couronne) est une ionisation locale de l'air autour
d'un conducteur à fort champ de surface ($E_{\text{surface}} > E_{\text{critique}}$) sans
qu'il y ait claquage complet.

\begin{figure}[H]
\centering
\begin{tikzpicture}[scale=1.1]
  % Dark background circle
  \fill[darkblue!90] (0,0) circle (3.5);
  % Corona glow layers
  \foreach \r/\op in {0.5/0.7, 0.8/0.5, 1.1/0.4, 1.5/0.3, 2.0/0.2, 2.6/0.1}{
    \draw[lightblue!\the\numexpr int(90*\op)\relax, opacity=\op,
          line width=\r*2pt] (0,0) circle (\r);
  }
  % Conductor
  \fill[white!70] (0,0) circle (0.3);
  \node[font=\footnotesize\bfseries, darkblue] at (0,0) {HT};
  % Field lines
  \foreach \a in {0,30,...,330}{
    \draw[-Stealth, accent!80, opacity=0.8, line width=1pt]
      ({0.35*cos(\a)},{0.35*sin(\a)}) -- ({2.5*cos(\a)},{2.5*sin(\a)});
  }
  % Ion markers
  \foreach \a/\r in {15/0.7, 60/0.9, 105/0.8, 160/1.0, 200/0.75, 250/0.85,
                     295/0.95, 340/0.7}{
    \node[font=\tiny, yellow!90!orange] at ({(\r)*cos(\a)},{(\r)*sin(\a)}) {$\star$};
  }
  % Labels
  \node[white, font=\small, anchor=west] at (2.7, 1.5) {Zone d'ionisation};
  \node[accent, font=\small, anchor=west] at (2.7, 0.8) {Lignes de champ $\vec{E}$};
  \node[yellow!80, font=\small, anchor=west] at (2.7, 0.1) {$\star$ Ions/électrons};
\end{tikzpicture}
\caption{Effet corona autour d'un conducteur cylindrique sous haute tension.
La zone bleue lumineuse représente la région ionisée (plasma partiel).}
\end{figure}

Pour un conducteur cylindrique de rayon $a$, le corona commence quand :
\begin{equation}
  E(r=a) = \frac{U}{a\,\ln(R/a)} > E_{\text{critique}} \approx 2.1\,\si{MV/m}
  \left(1 + \frac{0.03}{\sqrt{a}}\right) \quad \text{(formule de Peek)}
\end{equation}

\subsection{Conséquences de l'Effet Corona}

\begin{center}
\begin{tabularx}{\textwidth}{lXl}
\toprule
\rowcolor{darkblue}
\textcolor{white}{\textbf{Effet}} &
\textcolor{white}{\textbf{Description}} &
\textcolor{white}{\textbf{Ordre de grandeur}} \\
\midrule
Pertes d'énergie & Consommation parasite permanente & 0.5--5\% de la puissance \\
\rowcolor{lightgray}
Bruit audible & Sifflement/crépitement (100 Hz) & 40--65 dB(A) à 30 m \\
Interférences EM & Perturbations radio AM (0.5--30 MHz) & Jusqu'à 60 dB$\mu$V/m \\
\rowcolor{lightgray}
Production O$_3$ & Corrosion métaux, dégradation caoutchouc & 10--100 ng/J \\
Érosion conducteur & Piqûres de surface & Durée vie réduite \\
\bottomrule
\end{tabularx}
\end{center}

\section{Décharges Partielles (DP)}

\subsection{Définition et Modèle}

Les \textbf{décharges partielles} sont des ruptures locales d'isolation dans des
cavités gazeuses incluses dans l'isolation solide ou liquide. Elles ne traversent
pas complètement l'isolant mais l'érodent progressivement.

\begin{figure}[H]
\centering
\begin{tikzpicture}[scale=1.2]
  % Main dielectric
  \fill[accent!15] (0,0) rectangle (8,3);
  \draw[midblue, thick] (0,0) rectangle (8,3);
  % Electrodes
  \fill[darkblue!80] (0,3) rectangle (8,3.4);
  \fill[darkblue!80] (0,-0.4) rectangle (8,0);
  \node[white, font=\small\bfseries] at (4, 3.2) {Électrode HT};
  \node[white, font=\small\bfseries] at (4, -0.2) {Électrode GND};
  % Void (cavity)
  \fill[lightblue!40] (3.5,1.2) ellipse (1.2 and 0.5);
  \draw[hvred, thick] (3.5,1.2) ellipse (1.2 and 0.5);
  \node[hvred, font=\small\bfseries] at (3.5,1.2) {Cavité (gaz)};
  % Discharge symbol in void
  \draw[hvred, line width=2pt] (3.0,1.2) -- (3.3,1.5) -- (3.6,1.0) -- (3.9,1.4);
  \node[hvred, font=\scriptsize, above] at (3.5,1.7) {DP};
  % Equivalent circuit
  \node[font=\small, textdark] at (6.5,1.5) {$C_{\rm void}$, $C_{\rm series}$, $C_{\rm parallel}$};
  % Labels
  \node[textmed, font=\small] at (1.5,2.2) {Diélectrique};
  \node[textmed, font=\small] at (1.5,0.7) {Solide/Liquide};
\end{tikzpicture}
\caption{Modèle d'une décharge partielle dans une cavité gazeuse incluse
dans un diélectrique solide sous tension alternative.}
\end{figure}

La charge de DP est exprimée en \textbf{picocoulombs (pC)} :
\begin{equation}
  Q_{\rm DP} = C_{\rm void} \cdot \Delta V_{\rm cav} \quad [\si{pC}]
\end{equation}

\begin{warnbox}{Impact des DP sur la durée de vie}
Chaque décharge partielle dépose de l'énergie ($W = Q^2/(2C)$) qui érode chimiquement
et thermiquement l'isolant. Sous DP continues à 50 Hz :
\begin{itemize}
  \item Câble XLPE 150 kV : durée de vie réduite de 40 ans à quelques années si DP $>$ 50 pC.
  \item Transformateur : dégradation du papier accélérée, production de gaz (H$_2$, CH$_4$).
  \item Règle d'or : maintenir les DP en dessous des seuils IEC 60270 (voir tableau).
\end{itemize}
\end{warnbox}

%=============================================================================
\chapter{Surtensions}
%=============================================================================

\section{Classification des Surtensions}

\begin{center}
\begin{tabular}{llll}
\toprule
\rowcolor{darkblue}
\textcolor{white}{\textbf{Type}} &
\textcolor{white}{\textbf{Durée}} &
\textcolor{white}{\textbf{Amplitude}} &
\textcolor{white}{\textbf{Origine}} \\
\midrule
Temporaire (TOV)      & ms à min  & 1.2--1.5 $U_n$ & Défauts réseau, Ferranti \\
\rowcolor{lightgray}
Manœuvre (SOV)        & 0.1--20 ms & 2--4 $U_n$     & Commutation disjoncteurs \\
Foudre (LI)           & 1--100 µs  & 5--20 $U_n$    & Impact ou induction \\
\rowcolor{lightgray}
VFTO (GIS)            & 10 ns      & 2--2.5 $U_n$   & Opérations dans GIS \\
\bottomrule
\end{tabular}
\end{center}

\section{Surtensions Atmosphériques — Foudre}

\subsection{Physique du Coup de Foudre}

\begin{enumerate}
  \item Séparation des charges dans le cumulonimbus par processus de précipitation.
  \item Le champ au sol atteint 3--10 kV/m $\rightarrow$ ionisation descendante.
  \item Jonction traceur descendant / traceur ascendant à 100--300 m du sol.
  \item Arc de retour lumineux (foudre visible) : courant montant à $\approx 10^8\,\si{m/s}$.
  \item Courant crête : typiquement 30 kA (médiane), max. 300 kA (cas extrêmes).
\end{enumerate}

\subsection{Forme d'Onde Normalisée IEC 60060-1}

\begin{figure}[H]
\centering
\begin{tikzpicture}
\begin{axis}[
  width=0.85\textwidth, height=6cm,
  xlabel={Temps ($\mu$s)},
  ylabel={Tension normalisée $u/U_{\rm peak}$},
  xmin=-5, xmax=200,
  ymin=-0.1, ymax=1.2,
  grid=both, grid style={dotted, gray!40},
  axis background/.style={fill=lightgray!20},
  tick label style={font=\small},
  label style={font=\small},
  title={Onde de choc foudre normalisée 1.2/50 $\mu$s},
  title style={color=darkblue, font=\normalsize\bfseries},
]
% Lightning impulse shape approximation
\addplot[hvred, thick, line width=2pt, smooth, samples=300, domain=0:200]
  {max(0, (1 - exp(-x/1.2)) * exp(-x/50))/(0.368)};

% 90% and 30% lines
\addplot[dashed, midblue, opacity=0.7] coordinates {(-5,0.9) (200,0.9)}
  node[right, font=\scriptsize]{90\%};
\addplot[dashed, midblue, opacity=0.7] coordinates {(-5,0.3) (200,0.3)}
  node[right, font=\scriptsize]{30\%};
\addplot[dashed, accent, opacity=0.8] coordinates {(-5,0.5) (200,0.5)}
  node[right, font=\scriptsize]{50\% ($T_2$)};
\addplot[dashed, hvgreen, opacity=0.7] coordinates {(-5,1.0) (200,1.0)}
  node[right, font=\scriptsize]{$U_{\rm peak}$};

% Annotations
\draw[->, hvred, thick] (axis cs:1.2, 0) -- (axis cs:1.2, 0.9);
\node[hvred, font=\scriptsize, above] at (axis cs:5, 0.93) {$T_1=1.2\,\mu$s};
\draw[->, hvgreen!70!black, thick] (axis cs:50, 0) -- (axis cs:50, 0.5);
\node[accent, font=\scriptsize] at (axis cs:70, 0.53) {$T_2=50\,\mu$s};

\end{axis}
\end{tikzpicture}
\caption{Onde de choc foudre 1.2/50 µs selon IEC 60060-1. $T_1 = 1.2\,\mu$s (temps de
front), $T_2 = 50\,\mu$s (temps de queue à mi-crête).}
\end{figure}

\section{Surtensions de Manœuvre}

\begin{figure}[H]
\centering
\begin{tikzpicture}
\begin{axis}[
  width=0.85\textwidth, height=6cm,
  xlabel={Temps ($\mu$s)},
  ylabel={Tension normalisée},
  xmin=-50, xmax=8000,
  ymin=-0.1, ymax=1.25,
  grid=both, grid style={dotted, gray!40},
  axis background/.style={fill=lightgray!20},
  tick label style={font=\small},
  label style={font=\small},
  title={Onde de manœuvre normalisée 250/2500 $\mu$s},
  title style={color=darkblue, font=\normalsize\bfseries},
]
\addplot[midblue, thick, line width=2pt, smooth, samples=300, domain=0:8000]
  {(1 - exp(-x/250)) * exp(-x/2500) / 0.7829};

\addplot[dashed, accent, opacity=0.8] coordinates {(-50,1.0) (8000,1.0)}
  node[right, font=\scriptsize]{Crête};
\addplot[dashed, hvgreen, opacity=0.7] coordinates {(-50,0.5) (8000,0.5)}
  node[right, font=\scriptsize]{50\%};

\draw[->, midblue!80, thick] (axis cs:250, 0) -- (axis cs:250, 0.632);
\node[midblue, font=\scriptsize, above right] at (axis cs:250, 0.65)
  {$T_1=250\,\mu$s};
\draw[->, hvgreen!80, thick] (axis cs:2500, 0) -- (axis cs:2500, 0.5);
\node[hvgreen!70!black, font=\scriptsize, above right] at (axis cs:2500, 0.52)
  {$T_2=2500\,\mu$s};
\end{axis}
\end{tikzpicture}
\caption{Onde de manœuvre 250/2500 µs selon IEC 60060-1. Beaucoup plus lente que la
foudre, elle sollicite davantage l'isolation interne des équipements.}
\end{figure}

\section{Protection contre les Surtensions}

\subsection{Parafoudres à Oxyde Métallique (MOV)}

Les \textbf{parafoudres ZnO} (zinc oxide) sont les dispositifs de protection standard.
Leur caractéristique courant-tension est fortement non linéaire :
\begin{equation}
  I = k \cdot V^{\,\alpha} \quad \text{avec } \alpha = 25\text{--}100
\end{equation}

En service normal ($V < V_{\text{ref}}$) : $I < 1\,\si{mA}$ (quasi-isolant).\\
En surtension ($V > V_{\text{ref}}$) : limitation à $V_{\text{protection}} = V_{\text{ref}} \times (1 + \varepsilon)$.

\begin{retenir}{Sélection d'un parafoudre}
Le parafoudre est caractérisé par :
\begin{itemize}
  \item $U_c$ : tension continue de service maximale
  \item $U_r$ : tension assignée (tenue pendant 10 s)
  \item $U_p$ : niveau de protection (à $I_n$ nominal)
  \item $I_n$ : courant de choc nominal (typiquement 5--20 kA)
\end{itemize}
Critère de coordination : $U_p \times 1.15 \leq \text{NIL}$ de l'équipement protégé.
\end{retenir}

%=============================================================================
\chapter{Génération des Hautes Tensions}
%=============================================================================

\section{Génération Haute Tension Alternative (AC)}

\subsection{Transformateur de Test}

Le \textbf{transformateur de test} est le moyen le plus courant de générer des tensions
alternatives jusqu'à quelques MV. Pour des tensions très élevées, on utilise des montages
en \textbf{cascade} :

\begin{figure}[H]
\centering
\begin{tikzpicture}[scale=0.9,
  component/.style={draw=midblue, fill=midblue!10, rounded corners, minimum width=2cm, minimum height=1.2cm}
]
  % Stage 1
  \node[component, align=center] (T1) at (0,0) {\shortstack{\textbf{T$_1$}\\[2pt]\footnotesize $V_1 = 100\,\si{kV}$}};
  % Stage 2
  \node[component, align=center] (T2) at (3.5,1.5) {\shortstack{\textbf{T$_2$}\\[2pt]\footnotesize $V_2 = 100\,\si{kV}$}};
  % Stage 3
  \node[component, align=center] (T3) at (7,3.0) {\shortstack{\textbf{T$_3$}\\[2pt]\footnotesize $V_3 = 100\,\si{kV}$}};
  % Connections
  \draw[hvred, thick, -Stealth] (T1.north) -- ++(0,0.8) -| (T2.south west);
  \draw[hvred, thick, -Stealth] (T2.north) -- ++(0,0.8) -| (T3.south west);
  % Output
  \draw[hvred, very thick, -Stealth] (T3.east) -- ++(1.5,0)
    node[right, hvred, font=\bfseries\normalsize]{$V_{\rm out} = 300\,\si{kV}$};
  % Ground
  \draw[darkblue, thick] (T1.south) -- ++(0,-0.5) node[below, font=\small]{$\perp$ GND};
  % Excitation
  \draw[accent, thick, -Stealth] (-2,0) -- (T1.west)
    node[midway, above, font=\small, accent]{$V_{\rm in} = 10\,\si{kV}$};
  % Labels
  \node[textmed, font=\footnotesize] at (3.5,-0.8) {Montage en cascade (3 étages)};
\end{tikzpicture}
\caption{Transformateurs de test en cascade. Chaque étage isole uniquement sa propre
tension, simplifiant la conception de l'isolation à la masse.}
\end{figure}

\subsection{Résonance série — Test des Câbles}

Pour les câbles longs (forte capacité), la puissance réactive est compensée par une
inductance accordée en résonance :
\begin{equation}
  L_{\rm res} = \frac{1}{\omega^2 C_{\rm câble}}, \quad
  Q = \frac{V_{\rm câble}}{V_{\rm source}} = \frac{\omega L_{\rm res}}{R}
\end{equation}
Le facteur de qualité $Q$ peut atteindre 50--100, permettant de générer 100 kV
avec une alimentation de seulement 1--2 kV. Très économique pour les essais de câbles.

\section{Génération Haute Tension Continue (DC)}

\subsection{Multiplicateur de Cockroft-Walton}

Le \textbf{multiplicateur de Cockroft-Walton} (1932) génère une tension DC
très élevée à partir d'une source AC par rectification et multiplication :

\begin{equation}
  V_{\rm DC} = 2n\,V_{\rm AC,peak} - \frac{n(n+1)(2n+1)}{6f C}\,I_{\rm charge}
\end{equation}

où $n$ = nombre d'étages, $f$ = fréquence, $C$ = capacité des condensateurs.

\begin{figure}[H]
\centering
\begin{tikzpicture}[
  scale=1.0,
  diode/.style={draw=hvred, fill=hvred!20, isosceles triangle, isosceles triangle apex angle=60,
                minimum size=0.5cm},
]
  % Simplified 3-stage CW multiplier
  \def\Vac{0}; % just for labels

  % Nodes
  \foreach \n/\x/\y in {A/0/0, B/0/2, C/2/0, D/2/2, E/4/0, F/4/2, G/6/2}{
    \coordinate (\n) at (\x, \y);
  }
  % Capacitors (vertical)
  \foreach \cx/\cy/\lab in {0/1/C_1, 2/1/C_2, 4/1/C_3}{
    \draw[midblue, thick] (\cx,{\cy-0.3}) -- (\cx,{\cy-0.1});
    \draw[midblue, thick] ({\cx-0.3},{\cy-0.1}) -- ({\cx+0.3},{\cy-0.1});
    \draw[midblue, thick] ({\cx-0.3},{\cy+0.1}) -- ({\cx+0.3},{\cy+0.1});
    \draw[midblue, thick] (\cx,{\cy+0.1}) -- (\cx,{\cy+0.3});
    \node[midblue, font=\scriptsize, right] at ({\cx+0.35},\cy) {$\lab$};
  }
  % Diodes (horizontal)
  \foreach \dx/\dy in {1/0, 3/0, 5/0, 1/2, 3/2, 5/2}{
    \draw[hvred, thick, -Stealth] (\dx,\dy) -- ({\dx+0.6},\dy);
  }
  % Horizontal wires
  \draw[darkblue, thick] (0,0) -- (6,0);
  \draw[darkblue, thick] (0,2) -- (6,2);
  % AC source
  \draw[accent, thick] (-1,1) circle (0.5);
  \node[accent, font=\small] at (-1,1) {$\sim$};
  \draw[accent] (-0.5,1) -- (0,1);
  \draw[accent] (0,1) -- (0,0); \draw[accent] (0,1) -- (0,2);
  % Output
  \draw[hvred, very thick] (6,0) -- (6.5,0);
  \draw[hvred, very thick] (6,2) -- (6.5,2);
  \node[hvred, font=\normalsize\bfseries, right] at (6.5,1) {$+6\,V_{\rm peak}$};
  % GND
  \draw[darkblue] (0,0) -- (0,-0.5) node[below, font=\small]{GND};
\end{tikzpicture}
\caption{Multiplicateur de Cockroft-Walton (3 étages). Chaque étage double la tension :
$V_{\rm DC} = 2n\,V_{\rm AC,peak}$.}
\end{figure}

\section{Génération d'Impulsions — Générateur de Marx}

\subsection{Principe de Fonctionnement}

Le \textbf{générateur de Marx} (1923) est le circuit standard mondial pour les essais
d'impulsions de tension simulant la foudre. Son principe fondamental est simple et
élégant :

\begin{itemize}
  \item \textbf{Phase de charge :} Tous les condensateurs $C_1 \cdots C_n$ se chargent
    \textit{en parallèle} à la tension $V_c$ via les résistances $R_{\rm charge}$.
  \item \textbf{Déclenchement :} L'amorçage du premier éclateur en déclenche d'autres en
    cascade, reconnectant les condensateurs \textit{en série}.
  \item \textbf{Phase de décharge :} Tension de sortie $V_{\rm out} = n \cdot V_c$.
  \item \textbf{Mise en forme :} La résistance $R_1$ façonne le front, $R_2$ la queue.
\end{itemize}

\begin{figure}[H]
\centering
\begin{tikzpicture}[scale=1.0,
  cap/.style={draw=midblue, fill=midblue!10, rectangle, minimum width=0.3cm, minimum height=0.8cm},
  gap/.style={draw=hvred, fill=hvred!15, circle, minimum size=0.5cm, inner sep=1pt},
]
  % 4-stage Marx generator
  \def\stg{4}
  \foreach \i in {1,...,\stg}{
    \pgfmathsetmacro\x{(\i-1)*3.2}
    % Capacitor
    \fill[midblue!15] ({\x+0.5},0) rectangle ({\x+0.9}, 3.5);
    \draw[midblue, thick] ({\x+0.7}, 0.3) -- ({\x+0.7}, 1.5);
    \draw[midblue, thick] ({\x+0.3}, 1.5) -- ({\x+1.1}, 1.5);
    \draw[midblue, thick] ({\x+0.3}, 1.7) -- ({\x+1.1}, 1.7);
    \draw[midblue, thick] ({\x+0.7}, 1.7) -- ({\x+0.7}, 3.2);
    \node[midblue, font=\footnotesize] at ({\x+0.7},-0.4) {$C_{\i}$};
    \node[midblue, font=\scriptsize] at ({\x+0.7},-0.8) {$V_c=100\,\si{kV}$};
    % Spark gap
    \ifnum\i<\stg
      \fill[hvred!20] ({\x+1.8}, 1.6) circle (0.3);
      \draw[hvred, thick] ({\x+1.8}, 1.6) circle (0.3);
      \node[hvred, font=\tiny\bfseries] at ({\x+1.8},1.6) {EG};
      % Charging resistors (top)
      \draw[accent, thick, decorate, decoration={zigzag,amplitude=2pt,segment length=4pt}]
        ({\x+0.7},3.2) -- ({\x+3.2+0.7},3.2);
      \node[accent, font=\tiny, above] at ({\x+2.0},3.25) {$R_c$};
    \fi
  }
  % Output waveshaping
  \pgfmathsetmacro\xout{(\stg)*3.2}
  \draw[hvred, very thick, -Stealth] ({\xout+0.0},1.6) -- ({\xout+1.5},1.6)
    node[right, hvred, font=\bfseries]
      {$V_{\rm out}=\stg\times100 = 400\,\si{kV}$};
  % R1 (front)
  \draw[hvred, thick, decorate, decoration={zigzag,amplitude=2pt,segment length=4pt}]
    ({\xout+0.0},1.6) -- ({\xout+0.0},3.6);
  \node[hvred, font=\small, right] at ({\xout+0.1},2.6) {$R_1$ (front)};
  % R2 (tail)
  \draw[hvred, thick, decorate, decoration={zigzag,amplitude=2pt,segment length=4pt}]
    ({\xout+0.0},1.6) -- ({\xout+0.0},-0.4);
  \node[hvred, font=\small, right] at ({\xout+0.1},0.6) {$R_2$ (queue)};
  % GND
  \draw[darkblue, thick] (0.7,0.3) -- (0.7,-0.2) node[below, font=\small]{GND};
\end{tikzpicture}
\caption{Générateur de Marx à 4 étages. Chaque condensateur est chargé à 100 kV.
Au déclenchement, ils se mettent en série : $V_{\rm out} = 4 \times 100 = 400\,\si{kV}$.}
\end{figure}

\begin{exemple}{Calcul d'un générateur de Marx 10 étages — 1 MV}
\textbf{Données :} $n = 10$ étages, $V_c = 100\,\si{kV}$ par condensateur, $C = 0.5\,\si{\mu F}$\\[6pt]
\textbf{Tension de sortie :}
\[ V_{\rm out} = n \times V_c = 10 \times 100 = 1\,000\,\si{kV} = \mathbf{1\,MV} \]
\textbf{Énergie stockée :}
\[ W = \frac{1}{2} n C V_c^2 = \frac{1}{2}\times10\times0.5\times10^{-6}\times(10^5)^2 = 25\,\si{kJ} \]
\textbf{Résistances de mise en forme :}
\begin{align*}
  R_1 &= \frac{T_1}{0.7 \times n\,C} = \frac{1.2\times10^{-6}}{0.7\times10\times0.5\times10^{-6}} \approx 0.34\,\Omega \\
  R_2 &= \frac{T_2}{0.7 \times C_{\rm load}} \quad \text{(dépend de la charge)}
\end{align*}
\end{exemple}

%=============================================================================
\chapter{Mesure des Hautes Tensions}
%=============================================================================

\section{Éclateur à Sphères (Sphere Gap)}

L'\textbf{éclateur à sphères} (IEC 60052) est la méthode de référence absolue.
La tension de claquage entre deux sphères identiques de diamètre $D$ séparées d'une
distance $d$ est donnée par des tables normalisées.

\begin{figure}[H]
\centering
\begin{tikzpicture}[scale=1.1]
  % HV sphere
  \fill[midblue!40] (-3,0) circle (1.4);
  \draw[midblue, thick, line width=1.5pt] (-3,0) circle (1.4);
  \node[darkblue, font=\normalsize\bfseries] at (-3,0) {HT};
  \draw[hvred, very thick] (-3,2.5) -- (-3,1.4);
  \node[hvred, font=\small, above] at (-3, 2.6) {$V_{\rm HT}$};
  % GND sphere
  \fill[darkblue!40] (3,0) circle (1.4);
  \draw[darkblue, thick, line width=1.5pt] (3,0) circle (1.4);
  \node[white, font=\normalsize\bfseries] at (3,0) {GND};
  \draw[darkblue, thick] (3,-1.4) -- (3,-2.2);
  % Ground symbol
  \draw[darkblue, thick] (2.3,-2.2) -- (3.7,-2.2);
  \draw[darkblue, thick] (2.6,-2.5) -- (3.4,-2.5);
  \draw[darkblue, thick] (2.9,-2.8) -- (3.1,-2.8);
  % Gap
  \draw[<->, thick, accent] (-1.6,0) -- (1.6,0) node[midway, above, font=\bfseries\large]{$d$};
  % Spark (zigzag)
  \draw[yellow!80!orange, very thick, line width=2.5pt]
    (-1.5,0) -- (-1.0,0.3) -- (-0.5,-0.2) -- (0,0.25) -- (0.5,-0.2) -- (1.0,0.3) -- (1.5,0);
  \node[yellow!60!orange, font=\small, above] at (0, 0.5) {Étincelle};
  % Formula box
  \fill[boxbg] (4.5,-1.0) rectangle (7.5, 1.5);
  \draw[lightblue, thick] (4.5,-1.0) rectangle (7.5, 1.5);
  \node[darkblue, font=\small\bfseries] at (6.0, 1.1) {Correction atm.};
  \node[darkblue, font=\small] at (6.0, 0.55)
    {$\delta = \dfrac{p \times 293}{101.3 \times T}$};
  \node[darkblue, font=\small] at (6.0, -0.05) {$V = V_{\rm table} \times \delta$};
  \node[textmed, font=\scriptsize] at (6.0, -0.6) {$p$ [kPa], $T$ [K]};
  % Diameter label
  \draw[<->, thick, midblue] (-4.4,0) -- (-3,0) node[midway, below, font=\small]{$D/2$};
\end{tikzpicture}
\caption{Éclateur à sphères : dispositif de mesure de la valeur crête de la tension.
La table IEC 60052 donne $V_b = f(D, d)$ pour l'air sec à 20°C, 101.3 kPa.}
\end{figure}

\textbf{Exemple} : sphères $D = 25\,\si{cm}$, écartement $d = 5\,\si{cm}$, $p = 98\,\si{kPa}$, $T = 25°C$ :
\[
  V_{\rm table}(D=25\,\si{cm},\,d=5\,\si{cm}) = 93.2\,\si{kV} \quad \text{(IEC 60052 Table 2)}
\]
\[
  \delta = \frac{98 \times 293}{101.3 \times 298} = 0.953
  \qquad V_b = 93.2 \times 0.953 = \mathbf{88.8\,\si{kV}}
\]

\section{Diviseurs de Tension}

\subsection{Diviseur Résistif}

\begin{wrapfigure}{r}{0.42\textwidth}
\vspace{-8pt}
\begin{tikzpicture}[scale=1.05]
  % R1 (HV)
  \fill[midblue!15] (0.5,3.5) rectangle (2.5,6.5);
  \draw[midblue, thick, decorate,
        decoration={zigzag, amplitude=3pt, segment length=6pt}]
    (1.5,3.6) -- (1.5,6.4);
  \draw[midblue, thick] (1.5,6.4) -- (1.5,7.0);
  \draw[midblue, thick] (1.5,3.6) -- (1.5,2.8);
  \node[midblue, font=\bfseries, right] at (2.6,5.0) {$R_1$};
  \node[textmed, font=\scriptsize, right] at (2.6,4.5) {(HV)};
  % R2 (LV)
  \fill[hvgreen!15] (0.5,0.5) rectangle (2.5,2.6);
  \draw[hvgreen, thick, decorate,
        decoration={zigzag, amplitude=3pt, segment length=6pt}]
    (1.5,0.6) -- (1.5,2.5);
  \draw[hvgreen, thick] (1.5,2.5) -- (1.5,2.8);
  \draw[hvgreen, thick] (1.5,0.6) -- (1.5,0.0);
  \node[hvgreen, font=\bfseries, right] at (2.6,1.5) {$R_2$};
  \node[textmed, font=\scriptsize, right] at (2.6,1.0) {(LV)};
  % GND
  \draw[darkblue, thick] (0.8,0.0) -- (2.2,0.0);
  \draw[darkblue, thick] (1.0,-0.2) -- (2.0,-0.2);
  \draw[darkblue, thick] (1.2,-0.4) -- (1.8,-0.4);
  % HV input
  \draw[hvred, very thick, -Stealth] (1.5,8.0) -- (1.5,7.0);
  \node[hvred, font=\bfseries, above] at (1.5,8.1) {$V_{\rm HT}$};
  % Measurement point
  \fill[accent] (1.5,2.8) circle (0.12);
  \draw[accent, thick] (1.5,2.8) -- (4.0,2.8);
  \draw[accent, thick] (4.0,0.0) -- (4.0,2.8);
  % Voltmeter
  \draw[accent, thick] (3.5,0.0) -- (4.5,0.0) -- (4.5,2.5) -- (3.5,2.5) -- cycle;
  \node[accent, font=\bfseries\large] at (4.0,1.25) {V};
  \node[textmed, font=\scriptsize, right] at (4.6,1.25) {$V_m$};
  % Formula
  \node[darkblue, font=\normalsize] at (1.5,-1.1)
    {$V_{\rm HT} = V_m \dfrac{R_1+R_2}{R_2}$};
\end{tikzpicture}
\caption{Diviseur résistif HT}
\end{wrapfigure}

Le diviseur résistif est la méthode la plus courante pour la mesure permanente DC et AC.
Son rapport de division est :
\begin{equation}
  \boxed{k = \frac{V_{\rm HT}}{V_m} = \frac{R_1 + R_2}{R_2}}
\end{equation}

\textbf{Exemple} : $R_1 = 10\,\si{M\Omega}$, $R_2 = 10\,\si{k\Omega}$
\[
  k = \frac{10\,000 + 10}{10} = 1001 \approx 1000
\]
Pour $V_m = 100\,\si{V}$ : $V_{\rm HT} = 100\,\si{kV}$.

\medskip
\textbf{Limitations du diviseur résistif :}
\begin{itemize}
  \item Capacité parasite de $R_1$ dégrade la réponse au-delà de 1--10 kHz.
  \item Pour les impulsions HT, préférer le diviseur résistif-capacitif (RC).
  \item Consommation d'énergie : $P = V_{\rm HT}^2/(R_1+R_2)$.
\end{itemize}

\medskip

\subsection{Comparaison des Méthodes de Mesure HT}

\begin{center}
\begin{tabular}{lcccl}
\toprule
\rowcolor{darkblue}
\textcolor{white}{\textbf{Méthode}} &
\textcolor{white}{\textbf{DC}} &
\textcolor{white}{\textbf{AC}} &
\textcolor{white}{\textbf{Impulsions}} &
\textcolor{white}{\textbf{Précision}} \\
\midrule
Éclateur à sphères & -- & \checkmark & \checkmark & $\pm3\%$ \\
\rowcolor{lightgray}
Diviseur résistif   & \checkmark & \checkmark & -- & $\pm0.1\%$ \\
Diviseur capacitif  & -- & \checkmark & \checkmark & $\pm1\%$ \\
\rowcolor{lightgray}
Diviseur RC         & -- & \checkmark & \checkmark & $\pm1\%$ \\
TPC (transfo capac.)& -- & \checkmark & -- & $\pm0.2\%$ \\
\rowcolor{lightgray}
Pockels (électro-opt.)& \checkmark & \checkmark & \checkmark & $\pm0.5\%$ \\
\bottomrule
\end{tabular}
\end{center}

%=============================================================================
\chapter{Coordination de l'Isolation}
%=============================================================================

\section{Principe Général}

La \textbf{coordination de l'isolation} (IEC 60071) est le processus de sélection
des niveaux d'isolation des équipements d'un réseau électrique pour :
\begin{enumerate}
  \item Résister aux surtensions attendues en service normal.
  \item Permettre la protection par des limiteurs de surtension (parafoudres).
  \item Assurer la sécurité des personnes et l'intégrité des biens.
  \item Optimiser le coût global (isolation + protection + maintenance + indisponibilité).
\end{enumerate}

\section{Hiérarchie des Tensions Caractéristiques}

\begin{figure}[H]
\centering
\begin{tikzpicture}[scale=0.85]
\begin{axis}[
  ybar,
  symbolic x coords={
    {$U_n$}, {$U_m$}, {$U_p$ (parafoudre)},
    {NIL foudre}, {NIL manœuvre}
  },
  xtick=data,
  x tick label style={font=\small, rotate=20, anchor=east},
  ymin=0, ymax=1400,
  ylabel={Tension (kV)},
  ylabel style={font=\small},
  bar width=1.5cm,
  title={Coordination d'isolation — Réseau 400 kV},
  title style={color=darkblue, font=\bfseries\normalsize},
  axis background/.style={fill=lightgray!20},
  ymajorgrids=true, grid style={dotted, gray!40},
  width=0.92\textwidth, height=7cm,
]
\addplot[fill=hvgreen!60, draw=hvgreen!80] coordinates
  {({$U_n$},400) ({$U_m$},420) ({$U_p$ (parafoudre)},840)
   ({NIL foudre},1300) ({NIL manœuvre},1050)};
\end{axis}
\node[accent, font=\small] at (2.0,3.0) {$U_n=400$};
\node[midblue, font=\small] at (3.8,3.2) {$U_m=420$};
\node[hvred, font=\small] at (5.7,4.5) {$U_p=840$};
\end{tikzpicture}
\caption{Niveaux de tension pour la coordination d'isolation d'un réseau 400 kV.}
\end{figure}

\begin{formulebox}{Démarche de coordination — Méthode déterministe (IEC 60071-1)}
\begin{enumerate}
  \item Déterminer $U_m$ (tension max. d'équipement) : $U_m = 1.05\,U_n$
  \item Estimer la surtension de manœuvre maximale : $U_s = k_s \cdot U_m\sqrt{2/3}$
    avec $k_s = 3$ pour les réseaux non compensés.
  \item Choisir le parafoudre : $U_r \geq 1.25\,U_m/\sqrt{3}$
  \item Déterminer le niveau de protection : $U_p = f(U_r, I_n)$ (table fabricant)
  \item Calculer le NIL requis : $\text{NIL} \geq K_{cs} \cdot U_p$ ($K_{cs} = 1.15$ typique)
  \item Sélectionner le NIL normalisé supérieur dans la série IEC.
\end{enumerate}
\end{formulebox}

\begin{exemple}{Coordination d'isolation — Réseau 225 kV}
\begin{align*}
  U_n &= 225\,\si{kV} \\
  U_m &= 245\,\si{kV} \quad \text{(max. équipement IEC)} \\
  U_{m,\text{crête}} &= 245 \times \sqrt{2/3} = 200\,\si{kV} \\
  U_s &= 3 \times 200 = 600\,\si{kV} \quad \text{(surtension manœuvre)} \\
  U_r &\geq 1.25 \times 245/\sqrt{3} = 177\,\si{kV} \rightarrow \text{choisir } U_r = 192\,\si{kV} \\
  U_p &= 530\,\si{kV} \quad \text{(catalogue, } I_n = 10\,\si{kA)} \\
  \text{NIL} &\geq 1.15 \times 530 = 610\,\si{kV} \rightarrow \text{NIL normalisé} = \mathbf{650\,\si{kV}}
\end{align*}
\end{exemple}

%=============================================================================
\chapter{Applications Industrielles}
%=============================================================================

\section{Transformateurs de Puissance}

Les transformateurs de puissance (THT) sont les équipements les plus critiques et
les plus coûteux des réseaux électriques. Un transformateur 400/225 kV de 1000 MVA
représente un investissement de plusieurs millions d'euros et une durée de vie visée
de 40--50 ans.

\subsection{Construction et Isolation}

\begin{center}
\begin{tabularx}{\textwidth}{lX}
\toprule
\rowcolor{darkblue}
\textcolor{white}{\textbf{Partie}} &
\textcolor{white}{\textbf{Isolation et caractéristiques}} \\
\midrule
Enroulements HT/BT & Papier kraft imprégné huile (10--25 couches, $\varepsilon_r \approx 3.8$) \\
\rowcolor{lightgray}
Huile isolante & Huile naphténique ou ester, $E_{\rm claquage} > 70\,\si{kV/cm}$ \\
Traversées HT  & Condensateur céramique ou résine, NIL = 1300 kV (côté 400 kV) \\
\rowcolor{lightgray}
Noyau magnétique & Acier siliceux (3\% Si), mise à la terre en un seul point \\
Prise en charge   & Régleur en charge (OLTC), $\pm 10\,\%$ en 19 prises \\
\bottomrule
\end{tabularx}
\end{center}

\subsection{Essais Diélectriques Normalisés (IEC 60076-3)}

\begin{itemize}
  \item \textbf{Tenue à la tension appliquée (ACSD) :} $1.1 \times \sqrt{2} \times U_m/\sqrt{3}$
    pendant 60 s. Vérifie la tenue de l'isolation principale.
  \item \textbf{Tenue aux chocs de foudre (LI) :} 3 chocs positifs et 3 négatifs
    à l'amplitude NIL = 1300 kV pour un réseau 400 kV.
  \item \textbf{Mesure des décharges partielles :} $Q_{\rm max} < 500\,\si{pC}$ à
    $1.5\,U_m/\sqrt{3}$ pendant 30 min. Critère essentiel de santé de l'isolation.
  \item \textbf{Tenue aux chocs de manœuvre (SI) :} Obligatoire pour $U_m \geq 300\,\si{kV}$.
    Niveau : 1050 kV pour le réseau 400 kV.
\end{itemize}

\section{Lignes de Transport THT}

\begin{figure}[H]
\centering
\begin{tikzpicture}[scale=0.85]
  % Pylon
  \draw[darkblue, line width=2pt] (7,0) -- (7,9);
  \draw[darkblue, line width=1.5pt] (5.5,0) -- (7,2) -- (8.5,0);
  \draw[darkblue, line width=1.5pt] (4,3.5) -- (10,3.5);
  \draw[darkblue, line width=1.5pt] (4.5,6.0) -- (9.5,6.0);
  \draw[darkblue, line width=1.5pt] (5.2,8.0) -- (8.8,8.0);
  % Guard cable
  \draw[hvgreen, line width=1.5pt] (5.2,8.0) arc (180:0:1.8 and 0.3);
  \node[hvgreen, font=\tiny, above] at (7,8.35) {câble garde};
  % Insulator strings
  \foreach \x in {4, 10}{
    \foreach \dy in {0,...,4}{
      \draw[accent, line width=1.2pt] (\x, {3.5-\dy*0.22}) circle (0.08);
    }
    \draw[accent] (\x, 3.5) -- (\x, {3.5-4*0.22});
  }
  \foreach \x in {4.5, 9.5}{
    \foreach \dy in {0,...,4}{
      \draw[accent, line width=1.2pt] (\x, {6.0-\dy*0.22}) circle (0.07);
    }
    \draw[accent] (\x, 6.0) -- (\x, {6.0-4*0.22});
  }
  % Phase conductors (bundles)
  \foreach \x/\y in {4/2.4, 10/2.4, 4.5/4.9, 9.5/4.9}{
    \fill[midblue!60] (\x,\y) circle (0.18);
    \fill[midblue!30] ({\x+0.25},\y) circle (0.12);
    \fill[midblue!30] ({\x-0.25},\y) circle (0.12);
    \node[white, font=\tiny] at (\x,\y) {HT};
  }
  % Ground
  \draw[darkblue, very thick] (4,0) -- (10,0);
  % Labels
  \node[darkblue, font=\footnotesize, right] at (10.2,3.5) {$h_{\rm conduc.}$};
  \draw[<->, darkblue] (10.5,0) -- (10.5,3.5) node[midway,right,font=\tiny]{$\geq18\,\si{m}$};
  \node[accent, font=\footnotesize, left] at (3.5,3.0) {Chaîne isolateurs};
  \node[midblue, font=\footnotesize, left] at (3.3,2.4) {Faisceau 2 fils};
\end{tikzpicture}
\caption{Pylône de ligne 400 kV avec chaînes d'isolateurs, conducteurs en faisceau
et câble de garde supérieur.}
\end{figure}

\begin{retenir}{Caractéristiques des lignes THT françaises (RTE, 400 kV)}
\begin{itemize}
  \item Conducteurs : aluminium-acier (ACSR), section 570 mm$^2$, 2 conducteurs/phase
  \item Chaînes d'isolateurs : 22--24 disques verre par chaîne
  \item Distance d'isolement dans l'air : 2.5 m (phase-pylône), 6.5 m (phase-phase)
  \item Hauteur minimale au sol : 8 m (zone rurale), 10 m (voirie)
  \item Câble de garde : OPGW (fibre optique intégrée)
\end{itemize}
\end{retenir}

\section{Postes Blindés à Isolation Gazeuse (GIS)}

Les postes \textbf{GIS} (Gas-Insulated Switchgear) utilisent le SF$_6$ sous pression
(3--5 bar) comme isolant et extincteur. Ils représentent l'état de l'art en matière
de fiabilité et de compacité.

\begin{center}
\begin{tabularx}{\textwidth}{lXX}
\toprule
\rowcolor{darkblue}
\textcolor{white}{\textbf{Critère}} &
\textcolor{white}{\textbf{Poste conventionnel (air)}} &
\textcolor{white}{\textbf{Poste GIS}} \\
\midrule
Emprise au sol & $100\,\%$ (référence) & $5$--$10\,\%$ \\
\rowcolor{lightgray}
Entretien & Inspection tous 3--5 ans & Inspection tous 20--25 ans \\
Fiabilité (MTBF) & $\sim$50 ans (compos.) & $\sim$100 ans (compos.) \\
\rowcolor{lightgray}
Sensibilité météo & Forte (pluie, pollution) & Nulle (hermétique) \\
Coût d'installation & Faible & 3--5× plus élevé \\
\rowcolor{lightgray}
Bruit corona & Présent & Absent \\
\bottomrule
\end{tabularx}
\end{center}

%=============================================================================
\chapter{Normes et Standards}
%=============================================================================

\section{Normes IEC — Haute Tension}

\begin{longtable}{lp{7cm}p{4cm}}
\caption{Normes IEC principales pour la haute tension}\\
\toprule
\rowcolor{darkblue}
\textcolor{white}{\textbf{Norme}} &
\textcolor{white}{\textbf{Titre et domaine}} &
\textcolor{white}{\textbf{Édition}} \\
\midrule
\endfirsthead
\multicolumn{3}{c}{\small\itshape Suite du tableau~\thetable}\\
\toprule
\rowcolor{darkblue}
\textcolor{white}{\textbf{Norme}} &
\textcolor{white}{\textbf{Titre et domaine}} &
\textcolor{white}{\textbf{Édition}} \\
\midrule
\endhead
IEC 60060-1 & Techniques d'essais à haute tension — Mesures et définitions générales & 2010 \\
\rowcolor{lightgray}
IEC 60060-2 & Systèmes de mesure pour les essais HT & 2010 \\
IEC 60060-3 & Définitions et exigences pour les essais sur site & 2006 \\
\rowcolor{lightgray}
IEC 60071-1 & Coordination de l'isolation — Définitions, principes & 2019 \\
IEC 60071-2 & Coordination — Guide d'application & 2018 \\
\rowcolor{lightgray}
IEC 60076-3 & Transformateurs de puissance — Niveaux d'isolation & 2013 \\
IEC 60270   & Mesure des décharges partielles & 2015 \\
\rowcolor{lightgray}
IEC 60052   & Éclateurs à sphères — Tables de mesure & 2002 \\
IEC 60099-4 & Parafoudres — Parafoudres à oxyde métallique & 2014 \\
\rowcolor{lightgray}
IEC 60840   & Câbles HT $> 30\,\si{kV}$ — Essais & 2020 \\
IEC 62067   & Câbles THT $> 150\,\si{kV}$ — Essais & 2011 \\
\rowcolor{lightgray}
IEC 62271-203 & Appareillage HTA/HTB à isolation gazeuse (GIS) & 2022 \\
IEC 62271-100 & Disjoncteurs HT & 2021 \\
\rowcolor{lightgray}
IEC 60156   & Détermination de la rigidité diélectrique de l'huile & 2018 \\
\bottomrule
\end{longtable}

\section{Normes IEEE}

Les normes IEEE (principalement utilisées en Amérique du Nord) couvrent les mêmes
domaines avec quelques différences d'approche :

\begin{itemize}
  \item \textbf{IEEE Std 4} (2013) : Techniques d'essais HT — équivalent de IEC 60060.
  \item \textbf{IEEE Std 1313.1} (1996) : Coordination de l'isolation — Définitions.
  \item \textbf{IEEE Std 1313.2} (1999) : Coordination — Guide d'application.
  \item \textbf{IEEE C62.11} (2012) : Parafoudres ZnO pour systèmes AC.
  \item \textbf{IEEE C62.22} (2009) : Guide d'application des parafoudres.
  \item \textbf{IEEE C62.82.1} (2010) : Niveaux d'isolation pour systèmes $\leq 1000\,\si{kV}$.
\end{itemize}

\begin{infobox}{Convergence IEC--IEEE}
Les éditions récentes des normes IEC et IEEE tendent à s'harmoniser pour faciliter
les échanges internationaux d'équipements. Les valeurs numériques (NIL, tensions normalisées)
restent cependant légèrement différentes, ce qu'il faut toujours vérifier lors de projets
internationaux.
\end{infobox}

%=============================================================================
\chapter{Conclusion et Perspectives}
%=============================================================================

\section{Résumé des Concepts Clés}

Ce document a présenté les bases de l'ingénierie haute tension en couvrant :

\begin{enumerate}
  \item \textbf{Fondements électrostatiques :} Champ $\vec{E}$, potentiel $V$, loi de Gauss,
    matériaux diélectriques ($\varepsilon_r$, rigidité diélectrique).

  \item \textbf{Phénomènes de claquage :} Mécanisme de Townsend (gaz), streamers,
    claquage des liquides et solides, loi de Paschen, minimum de Paschen.

  \item \textbf{Effet corona et décharges partielles :} Mécanismes, conséquences industrielles,
    formule de Peek, mesure et limites normatives (IEC 60270).

  \item \textbf{Surtensions :} Foudre (1.2/50 µs), manœuvre (250/2500 µs),
    protection par parafoudres MOV (ZnO), câbles de garde.

  \item \textbf{Génération HT :} Transformateurs de test en cascade, multiplicateur
    Cockroft-Walton, générateur de Marx (principe charge-parallèle / décharge-série).

  \item \textbf{Mesure HT :} Éclateur à sphères (référence absolue), diviseurs
    résistifs/capacitifs/RC, méthodes électro-optiques (Pockels).

  \item \textbf{Coordination d'isolation :} Hiérarchie $U_n < U_m < U_s < U_p < \text{NIL}$,
    méthodes déterministe et probabiliste, normes IEC 60071.

  \item \textbf{Applications industrielles :} Transformateurs THT, lignes de transport,
    postes GIS (SF$_6$), problématique environnementale du SF$_6$.
\end{enumerate}

\section{Perspectives — Transition Énergétique et HT}

\subsection{HVDC et Réseaux Multiterminaux}

Le transport d'énergie en \textbf{courant continu haute tension (HVDC)} connaît
une croissance exponentielle avec la transition énergétique :

\begin{itemize}
  \item Liaisons terrestres $\pm$320 kV VSC pour les interconnexions européennes.
  \item Liaisons sous-marines $\pm$525 kV pour les parcs éoliens offshore
    (Dogger Bank UK, $\pm$500 MW par liaison).
  \item Vision ``SuperGrid'' : grille DC continentale reliant les sources
    renouvelables aux centres de consommation à l'échelle de l'Europe et de l'Afrique du Nord.
\end{itemize}

\subsection{Nouveaux Défis Techniques}

\begin{itemize}
  \item \textbf{Disjoncteurs DC HT :} Enjeu technologique majeur (pas d'annulation naturelle
    du courant). Solutions hybrides mécanique-électronique (ABB, GE, Siemens).
  \item \textbf{Isolation HVDC :} Le champ électrique dans les isolants DC dépend de la
    conductivité $\sigma$ (résistivité thermique) et non de $\varepsilon$ — comportement
    très différent du cas AC.
  \item \textbf{Remplacement du SF$_6$ :} Alternatives écologiques (mélanges g$^3$,
    CO$_2$/O$_2$, air sec) nécessitant une refonte de la conception des appareillages.
  \item \textbf{Supraconducteurs :} Câbles HTS (REBCO, BSCCO) pour le transport urbain
    haute densité de puissance sans pertes résistives.
  \item \textbf{Smart Grids :} Numérisation des postes (IEC 61850), capteurs IoT embarqués
    sur les équipements HT, jumeaux numériques, maintenance prédictive par IA.
\end{itemize}

\begin{retenir}{L'Ingénieur HT de Demain (2030--2050)}
Les compétences clés pour l'ingénieur HT de la prochaine décennie :
\begin{itemize}
  \item Maîtrise des convertisseurs VSC-HVDC (modélisation, protection, exploitation)
  \item Modélisation numérique avancée : FEM (COMSOL), transitoires (EMTP-ATP, PSCAD)
  \item Diagnostic et surveillance des équipements HT (DP, DGA, thermographie)
  \item Conception bas-carbone : zéro SF$_6$, matériaux biosourcés, économie circulaire
  \item Gestion des nouvelles sources de surtensions (onduleurs, FACTS, disjoncteurs DC)
\end{itemize}
\end{retenir}

\vspace{1cm}
\begin{center}
\textcolor{accent}{\rule{0.8\textwidth}{1.5pt}}\\[8pt]
{\large\textit{\textcolor{darkblue}{``Le génie électrique haute tension est la colonne vertébrale}}}\\
{\large\textit{\textcolor{darkblue}{de la transition énergétique mondiale.''}}}\\[8pt]
\textcolor{accent}{\rule{0.8\textwidth}{1.5pt}}
\end{center}

\vspace{1.5cm}
\noindent
\small
\textcolor{textmed}{
Ce document constitue une introduction structurée et pédagogique aux Techniques de
Haute Tension. Pour approfondir les sujets traités, les ouvrages suivants sont recommandés :
\begin{itemize}
  \item E.\ Kuffel, W.S.\ Zaengl \& J.\ Kuffel — \textit{High Voltage Engineering: Fundamentals}, Newnes, 2000.
  \item M.S.\ Naidu \& V.\ Kamaraju — \textit{High Voltage Engineering}, McGraw-Hill, 2013.
  \item A.\ Küchler — \textit{High Voltage Engineering: Fundamentals Technology — Applications}, Springer, 2018.
  \item C.L.\ Wadhwa — \textit{High Voltage Engineering}, New Age International, 2007.
\end{itemize}
}

\end{document}
